\begin{statementbox}
    \textbf{\textcolor{CStatement}{条件}}
    \begin{description}[style=nextline, leftmargin=2.8em, labelsep=0.5em, itemsep=0.35em]
        \item[\textbf{\textcolor{CStatement}{条件1（有公共端点）}}]
            平面直角坐标系中，$B(a,b)$、$P(x_1,y_1)$、$Q(x_2,y_2)$ 是三个点；两条线段 $BP$ 与 $BQ$ 都从点 $B$ 出发。
        \item[\textbf{\textcolor{CStatement}{条件2（垂直并且等长）}}]
            $\boxed{BP\perp BQ,\qquad BP=BQ>0}$。长度大于 $0$ 是为了排除点重合后无法确定方向的情况。
    \end{description}

\begin{center}
\begin{tikzpicture}[x=0.65cm,y=0.65cm,>=stealth,every node/.style={font=\small},line cap=round]

\draw[->,gray] (-3.6,0)--(4.8,0) node[right] {$x$};
\draw[->,gray] (0,-3.7)--(0,4.1) node[above] {$y$};
\draw[very thick,blue!75!black,->] (0,0)--(3,2);
\draw[very thick,teal!75!black,->] (0,0)--(-2,3);
\draw[very thick,orange!85!black,->] (0,0)--(2,-3);
\draw[teal!75!black,->,thick] (33.69:1.35) arc (33.69:123.69:1.35);
\draw[orange!85!black,->,thick] (33.69:1.85) arc (33.69:-56.31:1.85);
\draw (0.42,0.28)--(0.14,0.70)--(-0.28,0.42);
\fill (0,0) circle (2pt) node[below left] {$B(0,0)$};
\fill[blue!75!black] (3,2) circle (2pt) node[above right] {$P(3,2)$};
\fill[teal!75!black] (-2,3) circle (2pt) node[above left] {$Q_1(-2,3)$};
\fill[orange!85!black] (2,-3) circle (2pt) node[below right] {$Q_2(2,-3)$};
\node[teal!75!black,anchor=west] (ccw) at (4,3.3) {逆时针 $90^\circ$};
\node[orange!85!black,anchor=west] at (4,-2.1) {顺时针 $90^\circ$};
\node[anchor=west] at (4,1.6) {$BP=BQ_1=BQ_2$};
\node[anchor=west] at (4,0.7) {$=\sqrt{13}$};

\end{tikzpicture}
\end{center}

    \medskip
    \textbf{\textcolor{CStatement}{结论}}
    \begin{description}[style=nextline, leftmargin=2.8em, labelsep=0.5em, itemsep=0.45em]
        \item[\textbf{\textcolor{CStatement}{结论一（以原点为中心）}}]
            \textbf{\textcolor{CStatement}{直观描述：}} 水平、竖直两个坐标交换位置，其中一个变号。
            \[
            \boxed{\begin{array}{ll}
              \text{逆时针 }90^\circ:&(x,y)\longmapsto(-y,x),\\[3pt]
              \text{顺时针 }90^\circ:&(x,y)\longmapsto(y,-x).
            \end{array}}
            \]
        \item[\textbf{\textcolor{CStatement}{结论二（以任意点为中心）}}]
            设 $u=x_1-a,\ v=y_1-b$。以 $B(a,b)$ 为中心，把 $P$ 旋转 $90^\circ$ 后：
            \[
            \boxed{\begin{aligned}
              \text{逆时针：}\quad Q&=(a-v,\ b+u)=(a+b-y_1,\ b+x_1-a),\\
              \text{顺时针：}\quad Q&=(a+v,\ b-u)=(a+y_1-b,\ b+a-x_1).
            \end{aligned}}
            \]
            由于 $BP\perp BQ$ 且 $BP=BQ>0$，$Q$ 必是上面两点中的一个；反过来，任意一种坐标关系都保证上述两个几何条件。
        \item[\textbf{\textcolor{CStatement}{结论三（优先使用的消元关系）}}]
            不论顺、逆时针，恒有
            \[
              \boxed{|x_2-a|=|y_1-b|,\qquad |y_2-b|=|x_1-a|.}
            \]
            \textbf{特别有用：} 如果 $Q$ 位于直线 $x=a+d$，立即得到 $\boxed{|y_1-b|=|d|}$；若 $Q$ 位于 $y=b+d$，立即得到 $\boxed{|x_1-a|=|d|}$。
    \end{description}
    \begin{center}
\begin{tikzpicture}[x=.8cm,y=.8cm,>=stealth,every node/.style={font=\small}]
\draw[gray,->] (-4,0)--(2.6,0);
\draw[gray,dashed] (1,-.9)--(1,4.1) node[above] {$x=a+d$};
\draw[blue!75!black,very thick,->] (0,0)--(-3,1);
\draw[teal!75!black,very thick,->] (0,0)--(1,3);
\draw (-.36,.12)--(-.24,.48)--(.12,.36);
\draw[purple,dashed] (-3,0)--(-3,1) node[midway,left] {$v$};
\draw[purple,<->] (0,-.6)--(1,-.6) node[midway,below] {$d$};
\fill (0,0) circle (2pt) node[below left] {$B(a,b)$};
\node[above left] at (-3,1) {$P(x_1,y_1)$};
\node[right] at (1,3) {$Q(a+d,y_2)$};
\node[anchor=west] at (3,1.6) {横距固定，竖距随之固定};
\node[anchor=west] at (3,.4) {$|v|=|d|$};
\node at (0,-1.7) {图示 $d=v>0$；其余方向用绝对值统一表示。};
\end{tikzpicture}
\end{center}

    \medskip
    \textbf{\textcolor{CStatement}{严格区分}} 只有 $BP\perp BQ$、没有 $BP=BQ$ 时，只能判断方向垂直，不能直接用以上等长旋转公式。
\end{statementbox}
